\documentclass[11pt]{article}

\usepackage[T1]{fontenc}
\usepackage[utf8]{inputenc}
\usepackage{lmodern}
\usepackage{microtype}
\usepackage[margin=1in]{geometry}
\usepackage{amsmath,amssymb,amsthm,mathtools}
\usepackage{booktabs,tabularx,array}
\usepackage{enumitem}
\usepackage[hidelinks]{hyperref}

\newtheorem{theorem}{Theorem}[section]
\newtheorem{lemma}[theorem]{Lemma}
\newtheorem{proposition}[theorem]{Proposition}
\newtheorem{corollary}[theorem]{Corollary}
\theoremstyle{definition}
\newtheorem{definition}[theorem]{Definition}
\newtheorem{remark}[theorem]{Remark}

\newcommand{\F}{\mathbb F}
\newcommand{\Gone}{\mathbb G_1}
\newcommand{\Gtwo}{\mathbb G_2}
\newcommand{\GT}{\mathbb G_T}
\newcommand{\cC}{\mathcal C}
\newcommand{\cR}{\mathcal R}
\newcommand{\cL}{\mathcal L}
\newcommand{\cU}{\mathcal U}
\newcommand{\cV}{\mathcal V}
\newcommand{\Span}{\operatorname{span}}
\newcommand{\NAND}{\mathsf{NAND}}
\newcommand{\Setup}{\mathsf{Setup}}
\newcommand{\SimSetup}{\mathsf{SimSetup}}
\newcommand{\Prove}{\mathsf{Prove}}
\newcommand{\Vfy}{\mathsf{Vfy}}
\newcommand{\Sim}{\mathsf{Sim}}
\newcommand{\Ext}{\mathsf{Ext}}
\newcommand{\td}{\mathsf{td}}
\newcommand{\ProofObj}{\mathsf{proof}}
\newcommand{\cP}{\mathcal P}
\newcommand{\FSP}{\mathsf{FSP}}
\newcommand{\Enc}{\mathsf{Enc}}
\newcommand{\Dec}{\mathsf{Dec}}
\newcommand{\ProofMatrix}{\mathsf{ProofMatrix}}
\newcommand{\Test}{\mathsf{Test}}
\newcommand{\asg}{\operatorname{asg}}


\title{A Compact Split Non-interactive Linear Proof
for Boolean Circuit Satisfiability}
\author{Jens Groth}
\date{\today}

\begin{document}
\maketitle

\begin{abstract}
We construct a pairing-based zk-SNARK with one element from each source group in the generic Type-III bilinear-group model. Our construction does not use cryptographic hash functions, and the analysis does not rely on the random oracle model. The core of our construction is a non-interactive linear proof (NILP) for a sparse homogeneous linear system over signed Boolean values. The NILP is split, i.e., the field elements in the common reference string and the proof are partitioned into two separate tracks, which facilitates compilation into a Type-III pairing argument. Many algebraic pairing-based SNARKs can be viewed as compilations of split NILPs. This is the first split NILP with a single field element response on each side for a relation family whose induced language is NP-complete, and the first publicly verifiable preprocessing SNARK in the plain generic Type-III bilinear-group model whose proof contains one element from each source group. Verification is efficient; it requires work proportional to the public input. Proving is less efficient and uses a quadratic-size common reference string.
\end{abstract}

\section{Introduction}

Pairing-based succinct non-interactive arguments of knowledge (SNARKs) achieve small proofs by evaluating
polynomial identities in the exponent and checking quadratic relations between the proof elements with bilinear maps.
Groth's three-element construction arranges the
proof as two elements in the first source group and one in the second.  For
disclosure-free generic Type-III pairing-based arguments for hard decisional relation families, his lower bound requires proof elements from both source groups~\cite{Groth16}. This leaves a natural question: can such an argument have exactly one proof element in each source group?

Many pairing-based SNARKs can be viewed as instantiations of an information-theoretic abstraction called a non-interactive linear proof (NILP) defined over a finite field. Its common reference string generator samples field elements; without seeing the sampled CRS, the proof-matrix algorithm picks coefficients for linear combinations; the prover then applies those coefficients to the CRS; and the verifier checks the resulting proof using quadratic equations over the field.  In a split NILP, the common reference
string and the proof each have a left and a right part, each proof part is a
linear function only of its corresponding reference-string half, and the
verifier performs cross-bilinear tests.  A (disclosure-free) split NILP can be
compiled into a Type-III pairing argument in the generic bilinear-group model~\cite{Groth16}. The central problem we address is to construct a split NILP with a single field element response on each side for a relation family whose induced language is NP-complete.

Our construction has three algebraic layers. First, Boolean circuit
satisfiability is reduced to a sparse homogeneous linear system over signed
Boolean coordinates.  One residual coordinate per gate turns each NAND-gate constraint into a linear equation, leaving Booleanity as the only nonlinear condition on the coordinates.  Second, we introduce factorization span programs (FSPs).  An FSP specifies
two fixed linear spaces of polynomials and an input-linear family of public
targets.  Every target polynomial, defined by the public input, lies in the span of pairwise products from the two
spaces, but acceptance requires the stronger condition that the target polynomial admit a
single exact factorization into two factors, one from each space. Third, random evaluation of an FSP gives a split NILP with one field response on each side.

The main technical contribution is the FSP. The intuition behind it is to test Booleanity through unique polynomial factorization. To give an example, consider the polynomial $H(X)=X^2-1$ over a finite field $\F$. If $A(X)B(X)=H(X)$ with $A,B\in\F^{\leq 1}[X]$, then $A(X)$ must be of the form $v(X-b)$ for some $v\in\F^\times$ and $b\in\{-1,+1\}$. Using the convention $-1=\mathsf{false}$ and $+1=\mathsf{true}$, this means that $A(X)=v(X-b)$ represents a unique Boolean value $b$.

Generalizing this idea, we can represent many Boolean values $b_1,\ldots,b_n$ by a product of linear factors
\[
  A(X,Z)=\prod_{i=1}^n\bigl((Z-\omega_i)X-b_i\bigr),
  \qquad b_i\in\{-1,+1\},
\]
where $\omega_1,\ldots,\omega_n$ are distinct points in the field used as tags. Picking the complementary polynomial
\[
  B(X,Z)=\prod_{i=1}^n\bigl((Z-\omega_i)X+b_i\bigr),
\]
we have $A(X,Z)B(X,Z)=H(X,Z)$ for the fixed polynomial
\[
  H(X,Z)=\prod_{i=1}^n\bigl((Z-\omega_i)^2X^2-1\bigr).
\]
Unique factorization does not quite suffice to guarantee that $A(X,Z)$ is well formed. Even if $A(X,Z)$ has degree $n$ in $X$, an adversary might, for instance, include both $(Z-\omega_1)X-1$ and $(Z-\omega_1)X+1$ as factors, whereas we want to force the adversary to pick exactly one factor for each tag. We therefore require $A$ to belong to a specific linear space of polynomials $\cU$ whose leading-coefficient restriction ensures a unique choice of Boolean values. The restrictions defining $\cU$ can also impose linear constraints on $b_1,\ldots,b_n$. Finally, the full construction introduces a third indeterminate $Y$, and the target polynomial depends on the instance $\mathbf u$; we write it as $H_{\mathbf u}(X,Y,Z)$.

The split NILP over a finite field has perfect completeness, statistical knowledge soundness against split-affine strategies, and disclosure-freeness from the Schwartz--Zippel lemma. Multiplicative rerandomization gives perfect
simulation under a statistically close simulated CRS, and hence statistical
composable zero knowledge.  

Using the split NILP as the information-theoretic core, we can build a pairing-based preprocessing zk-SNARK in the generic Type-III bilinear-group model. The proof has one group element in each source group. While the proof has one fewer group element than Groth's construction, the CRS is quadratic
in the circuit size, so the construction is primarily of theoretical interest.
The main contribution is that this is the first
split NILP with split dimensions $(1,1)$ for a relation family whose induced
language is NP-complete, thus resolving a decade-old question~\cite[Section~1]{Groth16}.


\section{Boolean circuit satisfiability and homogeneous linear equations}
\label{sec:boolean-system}

Consider a Boolean circuit $\cC$, represented as a directed acyclic graph with
$G$ fan-in-two NAND gates and a single output.  Topologically order the gates
so that gate $G$ is the output gate, fix its output to true, and eliminate that
output-wire variable.  Let $W$ denote the number of remaining input and
internal wires.  Circuit satisfiability asks whether there exists a Boolean
assignment to the input wires such that the circuit outputs true.  This is
known to be an NP-complete problem.

For a fixed circuit $\cC$, we can designate $\ell$ of the input wires as public. We can then ask whether, for a given instance consisting of truth-value assignments to the public inputs, there exists a Boolean assignment to the remaining $W-\ell$ witness wires such that the circuit outputs true. This defines a finite relation $\cR_{\cC}$ between instances and witnesses, and a language $\cL_{\cC}$ of instances admitting a witness. NP-completeness applies to the resulting parameterized family: a polynomial-time constructible family of universal circuits can take the description of any suitably smaller circuit as part of its public input and evaluate that circuit on the remaining inputs~\cite{Valiant76}.

We can reduce a Boolean circuit $\cC$ to a sparse homogeneous linear system over signed Boolean values in a finite field $\F$. The reduction introduces one additional residual coordinate per gate, so
we set
\[
  N=W+G,
  \qquad
  n=N-\ell.
\]
Here $N$ is the number of nonconstant coordinates in the linear system, $\ell$ is the number of public input coordinates matching the circuit instance, and $n$ is the number of remaining coordinates. We will work over a finite field $\F$ of size
\[
  q=|\F|,
\]
and require $\operatorname{char}(\F)>3$ and $q>N$. Boolean values are
represented by signed coordinates in $\F$:
\[
  -1=\mathsf{false},
  \qquad
  +1=\mathsf{true}.
\]

Let $a_1,\ldots,a_W\in \{-1,+1\}$ represent truth values of the circuit wires.
After relabeling wires if necessary, the public input coordinates are
\[
  \mathbf u=(a_1,\ldots,a_\ell).
\]

Following~\cite{GOS06}, we can arithmetize the gates in the circuit. For each NAND gate $g$, with input coordinates $a_{i_g},a_{j_g}$ and output
coordinate $a_{k_g}$, introduce a fresh signed residual coordinate $a_{W+g}$ and
impose
\[
  -1+a_{i_g}+a_{j_g}+2a_{k_g}-a_{W+g}=0.
\]
For each gate, if and only if the gate is evaluated correctly, there exists a unique residual value $a_{W+g}\in\{-1,+1\}$ satisfying this equation. Specifically, the affine residual
\[
  r_g=-1+a_{i_g}+a_{j_g}+2a_{k_g}
\]
belongs to $\{-1,+1\}$ for a correct NAND evaluation and to
$\{-5,-3,3\}$ for an incorrect evaluation.  Setting $a_{W+g}=r_g$ therefore
gives the unique signed residual satisfying the gate equation exactly when the
gate is evaluated correctly.

It will be helpful to have a constant coordinate $a_0=1$ to homogenize the
linear system.  We represent the eliminated accepting output by $a_0$;
equivalently, set $k_G=0$.  For every other gate, $k_g$ is an ordinary
nonconstant wire index.
Thus each NAND gate becomes one homogeneous linear equation $-a_0+a_{i_g}+a_{j_g}+2a_{k_g}-a_{W+g}=0$ with at most five
nonzero coefficients. The word \emph{homogeneous} refers to the augmented
variables: the apparent constant term is the coefficient of $a_0$, whose value
is subsequently fixed to $1$.

We have focused on NAND gates since NAND is universal, but similar affine-map characterizations exist for all fan-in-two Boolean gates~\cite{DFGK14}.
Homogenizing each affine expression with $a_0$ and introducing a signed
residual coordinate gives the same kind of homogeneous gate equation, again
with at most five nonzero coefficients.

Each gate equation can be written as a column of a matrix $D\in\F^{(N+1)\times G}$, where the rows correspond to the augmented coordinates $(a_0,a_1,\ldots,a_W,a_{W+1},\ldots,a_N)$.  For the circuit reduction, the columns are automatically linearly independent: gate column $g$ has a nonzero entry in its dedicated residual row $W+g$, where every other gate column is zero in this row. Hence $\operatorname{rank}(D)=G$ for the concrete circuit matrix. The Boolean circuit satisfiability problem can therefore be reduced to finding a signed solution to a homogeneous linear system $(a_0,\ldots,a_N)D=0$ with $a_0=1$ and $a_1,\ldots,a_N\in \{-1,+1\}$, subject to the requirement that the first $\ell$ coordinates are fixed to the instance $\mathbf u$. The reduction is efficient and the resulting linear system has at most $5N$ nonzero entries, so $D$ is sparse.

Write
\[
  \mathbf u=(a_{1},\ldots,a_\ell), \qquad \mathbf w=(a_{\ell+1},\ldots,a_N),
  \qquad
  \overline{\mathbf a}(\mathbf u,\mathbf w)
  =(1,\mathbf u,\mathbf w).
\]
The Boolean homogeneous linear relation defined by $\F,\ell,G,W$, and $D\in \F^{(N+1)\times G}$, with $N=G+W$ and $n=N-\ell$, is
\[
  \cR
  =\left\{(\mathbf u,\mathbf w)\in
    \{-1,+1\}^{\ell}\times\{-1,+1\}^{n}:
    \overline{\mathbf a}(\mathbf u,\mathbf w)D=0\right\}.
\]
It induces the language
\[
  \cL
  =\left\{\mathbf u\in\{-1,+1\}^{\ell}:
    \exists\,\mathbf w\in\{-1,+1\}^{n}\text{ such that }
    (\mathbf u,\mathbf w)\in\cR\right\}.
\]
Thus $\mathbf u$ is the instance, the statement is the proposition
$\mathbf u\in\cL$, and $\mathbf w$ is a witness establishing that
statement. In the following, we will construct a proof system for this relation.



\section{Split non-interactive linear proofs}

Following Groth~\cite{Groth16}, we recall the NILP syntax in the row-vector
convention used throughout this paper.  A split NILP for a relation $\cR$
exposes the following algorithms.

\begin{definition}[Split NILP syntax]
Setup produces two public CRS row vectors and an associated trapdoor,
\[
  ((\boldsymbol{\sigma}_1,\boldsymbol{\sigma}_2),\td)
  \leftarrow\Setup(\cR),
  \qquad
  \boldsymbol{\sigma}_1\in\F^{d_1},\quad
  \boldsymbol{\sigma}_2\in\F^{d_2}.
\]
On an instance and witness $(\mathbf u,\mathbf w)\in\cR$, the randomized
proof-matrix algorithm computes
\[
  (P_1,P_2)\leftarrow\ProofMatrix(\cR,\mathbf u,\mathbf w),
  \qquad
  P_1\in\F^{k_1\times d_1},\quad
  P_2\in\F^{k_2\times d_2}.
\]
Crucially, $\ProofMatrix$ does not receive the sampled CRS.  The prover
algorithm does receive it and applies each proof matrix only to its
corresponding CRS half:
\[
  (\boldsymbol{\pi}_1,\boldsymbol{\pi}_2)
  \leftarrow
  \Prove(\cR,\boldsymbol{\sigma}_1,\boldsymbol{\sigma}_2,
         \mathbf u,\mathbf w),
  \qquad
  \boldsymbol{\pi}_1=\boldsymbol{\sigma}_1P_1^T,\qquad
  \boldsymbol{\pi}_2=\boldsymbol{\sigma}_2P_2^T.
\]
The deterministic test algorithm computes
\[
  (T_1,\ldots,T_\eta)\leftarrow\Test(\cR,\mathbf u),
  \qquad
  T_h\in\F^{(d_1+k_1)\times(d_2+k_2)},
\]
and $\Vfy$ first rejects unless $\mathbf u$ belongs to the public input
domain $\mathcal X$ of $\cR\subseteq\mathcal X\times\mathcal W$, and otherwise
accepts precisely when every cross-bilinear test vanishes:
\[
  (\boldsymbol{\sigma}_1,\boldsymbol{\pi}_1)
  T_h
  (\boldsymbol{\sigma}_2,\boldsymbol{\pi}_2)^T=0
  \qquad(1\le h\le\eta).
\]
The associated completeness, knowledge soundness, disclosure freeness, and zero-knowledge requirements are defined next.
\end{definition}

We use the following fixed-parameter security notions.  They quantify exact
error probabilities for the fixed relation and field under consideration and
therefore require neither a security parameter nor an asymptotic notion of
negligibility.  For zero knowledge, we separate closeness of the real and
simulated CRS distributions from perfect proof simulation under the simulated
CRS, following Groth's formulation of composable zero
knowledge~\cite{Groth06CZK}.

\begin{definition}[Fixed-parameter security]
\label{def:nilp-security}
A split NILP has the following properties.
\begin{enumerate}
  \item It has \emph{Perfect completeness} if, for every
  $(\mathbf u,\mathbf w)\in\cR$, an honestly generated proof is accepted with
  probability one over setup and prover randomness.

  \item It has \emph{split-affine statistical knowledge error at most
  $\varepsilon$} if there is an efficient extractor $\Ext$ such that the
  following holds.  Let any possibly randomized strategy, without seeing the
  sampled CRS, choose an instance $\mathbf u$ and proof matrices $P_1,P_2$ of
  the prescribed dimensions.  For an independently generated setup, set
  $\boldsymbol{\pi}_i=\boldsymbol{\sigma}_iP_i^T$.  Then
  \[
    \Pr\!\left[
      \Vfy(\cR,\boldsymbol{\sigma}_1,\boldsymbol{\sigma}_2,
           \mathbf u,(\boldsymbol{\pi}_1,\boldsymbol{\pi}_2))=1
      \ \land\
      \bigl(\mathbf u,\Ext(\cR,\mathbf u,P_1,P_2)\bigr)\notin\cR
    \right]
    \le\varepsilon.
  \]
  The word \emph{affine} allows constant terms; when each CRS half contains
  the coordinate $1$, they are represented by ordinary linear proof
  matrices.

  \item It has \emph{disclosure disagreement at most $\delta$} if, for every
  setup-independent matrix $T\in\F^{d_1\times d_2}$ and two independent
  setups $(\boldsymbol{\sigma}_1,\boldsymbol{\sigma}_2)$ and
  $(\boldsymbol{\sigma}'_1,\boldsymbol{\sigma}'_2)$,
  \[
    \Pr\!\left[
      \boldsymbol{\sigma}_1T\boldsymbol{\sigma}_2^T=0
      \text{ and }
      \boldsymbol{\sigma}'_1T(\boldsymbol{\sigma}'_2)^T\neq 0
    \right]
    \le\delta.
  \]
  We say the system has a $\delta$-\emph{disclosure-free} CRS.  

  \item It has \emph{statistical composable zero knowledge with error at most
  $\zeta$} if there are efficient algorithms $\SimSetup$ and $\Sim$ satisfying
  the following two properties.
  \begin{enumerate}[label=(\alph*)]
    \item The real and simulated public CRS distributions are statistically
    close.  Namely, if
    \[
      ((\boldsymbol{\sigma}_1,\boldsymbol{\sigma}_2),\td)
      \leftarrow\Setup(\cR)
      \quad\text{and}\quad
      ((\widehat{\boldsymbol{\sigma}}_1,
        \widehat{\boldsymbol{\sigma}}_2),\widehat{\td})
      \leftarrow\SimSetup(\cR),
    \]
    then
    \[
      \Delta\!\left(
        (\boldsymbol{\sigma}_1,\boldsymbol{\sigma}_2),
        (\widehat{\boldsymbol{\sigma}}_1,
         \widehat{\boldsymbol{\sigma}}_2)
      \right)\le\zeta.
    \]
    Only the public CRS distributions are compared; the setup trapdoors need
    not have the same distribution. In fact, in real usage $\td$ is not revealed and just represents the internal variables used by the setup algorithm.

    \item Proof simulation is perfect under the simulated CRS, even given the
    simulation trapdoor.  Give
    $(\widehat{\boldsymbol{\sigma}}_1,
      \widehat{\boldsymbol{\sigma}}_2,\widehat{\td})$ to any adaptive environment
    $\mathcal A$, which chooses an instance--witness pair
    $(\mathbf u,\mathbf w)$ as a function of the simulated CRS and trapdoor.
    If the pair is
    invalid, both experiments return $\bot$.  Otherwise, one experiment
    returns a fresh honest proof
    \[
      \Prove(\cR,\widehat{\boldsymbol{\sigma}}_1,
             \widehat{\boldsymbol{\sigma}}_2,\mathbf u,\mathbf w),
    \]
    while the other returns
    \[
      \Sim(\cR,\widehat{\boldsymbol{\sigma}}_1,
           \widehat{\boldsymbol{\sigma}}_2,\widehat{\td},\mathbf u),
    \]
    without giving $\mathbf w$ to $\Sim$.  The two full views, including the
    environment's randomness, the simulated CRS and trapdoor, its chosen pair,
    the reply, and its final output, must be identically distributed.
  \end{enumerate}
  Replacing the real CRS by the simulated CRS and then using perfect proof
  simulation shows that this definition implies the usual real-versus-
  simulated view distance of at most $\zeta$.  This is the statistical version
  of Groth's \emph{composable zero knowledge}, which is stronger than standard
  zero knowledge and implies unbounded adaptive zero knowledge under reuse of
  the same CRS~\cite{Groth06CZK}.  The property is \emph{perfect composable
  zero knowledge} when $\zeta=0$.
\end{enumerate}
\end{definition}

The pair $(k_1,k_2)$ is the split proof dimension.  Groth showed that a
disclosure-free generic Type-III pairing argument for a relation generator
with hard decisional problems must contain proof elements from both source
groups.  His construction has split proof dimension $(2,1)$.  We construct a
$(1,1)$-split NILP for a relation family whose induced language is NP-complete,
attaining the smallest split dimensions consistent with this lower bound.

\section{Factorization span programs}
\label{sec:fsp}

This section isolates the polynomial object underlying the construction.  It
plays the same organizational role that other polynomial objects such as QSP, QAP, SSP and SAP play in other
pairing-based proof systems: it is an algebraic relation, not yet
a proof system.  A generic random-evaluation compiler will later turn it into
a split NILP.

Unless stated otherwise, vectors are row vectors.  For an integer $d\ge0$,
write
\[
  \F^{\le d}[Z]=\{f(Z)\in\F[Z]:\deg f\le d\}.
\]
For a multivariate polynomial, $\deg_{\mathrm{tot}}$ denotes total degree.

\subsection{Definition}

Let
\[
  \cP=\F[X_1,\ldots,X_s]
\]
be a polynomial ring.  Let $\cU,\cV\subseteq\cP$ be finite-dimensional
$\F$-linear spaces with fixed row-vector bases
\[
  \mathbf U=(U_1,\ldots,U_{d_U}),
  \qquad
  \mathbf V=(V_1,\ldots,V_{d_V}).
\]
We require $1\in\cU\cap\cV$ and choose $U_1=V_1=1$.

\begin{definition}[Factorization span program]
\label{def:fsp}
An $\ell$-input \emph{factorization span program} (FSP) consists of the
spaces and bases above, target polynomials
\[
  H_0,H_1,\ldots,H_\ell\in\cP,
\]
and matrices
\[
  M_0,M_1,\ldots,M_\ell\in\F^{d_U\times d_V}
\]
such that
\begin{equation}
  \mathbf U(\mathbf X)M_i\mathbf V(\mathbf X)^T=H_i(\mathbf X)
  \qquad (0\le i\le\ell).
  \label{eq:fsp-target-representations}
\end{equation}
For a public input $\mathbf u=(u_1,\ldots,u_\ell)\in\F^\ell$, define
\begin{equation}
  H_{\mathbf u}(\mathbf X)
  =H_0(\mathbf X)+\sum_{i=1}^{\ell}u_iH_i(\mathbf X),
  \qquad
  M_{\mathbf u}=M_0+\sum_{i=1}^{\ell}u_iM_i.
  \label{eq:fsp-input-target}
\end{equation}
The program accepts $\mathbf u$ if there exist
$\mathsf A\in\cU$ and $\mathsf B\in\cV$ satisfying the formal identity
\begin{equation}
  \mathsf A(\mathbf X)\mathsf B(\mathbf X)
  =H_{\mathbf u}(\mathbf X).
  \label{eq:fsp-factorization}
\end{equation}
\end{definition}

The affine-linear dependence of $H_{\mathbf u}$ and $M_{\mathbf u}$ on the
public input is what we refer to as \emph{input linearity}.

\begin{remark}[Rank-one preimages]
Define the multiplication map
\[
  \mu:\F^{d_U\times d_V}\longrightarrow\cP,
  \qquad
  \mu(M)=\mathbf U M\mathbf V^T.
\]
The matrix $M_{\mathbf u}$ is an arbitrary preimage of the public target.
Writing
\[
  \mathsf A=\mathbf c_A\mathbf U^T,
  \qquad
  \mathsf B=\mathbf c_B\mathbf V^T,
\]
an accepting factorization gives the rank-one preimage
$\mathbf c_A^T\mathbf c_B\in\F^{d_U\times d_V}$.  Thus an FSP asks whether a
target known to lie in a product span $\cU \cV$ has a rank-one preimage under $\mu$.
\end{remark}

\subsection{Realizing an external relation}

The native witness of an FSP is a pair of polynomials.  To use the object for
an NP relation, we require efficient translations between ordinary witnesses
and polynomial factors.

\begin{definition}[FSP realization]
\label{def:fsp-realization}
Let $\cR\subseteq\mathcal X\times\mathcal W$, where
$\mathcal X\subseteq\F^\ell$.  An FSP \emph{realizes} $\cR$ if there are
polynomial-time algorithms $\Enc$ and $\Dec$ such that:
\begin{enumerate}[label=(\roman*)]
  \item if $(\mathbf u,\mathbf w)\in\cR$, then
  \[
    (\mathsf A,\mathsf B)\leftarrow\Enc(\mathbf u,\mathbf w)
  \]
  satisfies $\mathsf A\in\cU$, $\mathsf B\in\cV$, and
  $\mathsf A\mathsf B=H_{\mathbf u}$;
  \item if $\mathbf u\in\mathcal X$ and
  $\mathsf A\in\cU$, $\mathsf B\in\cV$ satisfy
  $\mathsf A\mathsf B=H_{\mathbf u}$, then
  \[
    \mathbf w\leftarrow\Dec(\mathbf u,\mathsf A,\mathsf B)
  \]
  satisfies $(\mathbf u,\mathbf w)\in\cR$.
\end{enumerate}
\end{definition}

For an FSP, define the degree parameter
\begin{equation}
  d_{\FSP}
  =\max\left\{
      \max_{i,j}\deg_{\mathrm{tot}}(U_iV_j),
      \max_{0\le k\le\ell}\deg_{\mathrm{tot}}H_k
    \right\}.
  \label{eq:fsp-degree}
\end{equation}

\subsection{Generic compiler to a split NILP}

The FSP abstraction makes the information-theoretic compiler immediate.

\begin{theorem}[FSP compiler]
\label{thm:fsp-compiler}
Suppose an $\ell$-input FSP over $\F$ realizes a relation $\cR$ and has
degree parameter $d_{\FSP}$.  Random evaluation gives a $(1,1)$-split NILP
with perfect completeness, split-affine statistical knowledge error at most
$d_{\FSP}/q$, and disclosure disagreement at most $d_{\FSP}/q$.
\end{theorem}

\begin{proof}
Sample independently and uniformly
\[
  \mathbf x=(x_1,\ldots,x_s)\leftarrow\F^s
\]
and let the two CRS halves be
\begin{equation}
  \boldsymbol{\sigma}_1=\mathbf U(\mathbf x)\in \F^{d_U},
  \qquad
  \boldsymbol{\sigma}_2=\mathbf V(\mathbf x)\in \F^{d_V}.
  \label{eq:generic-fsp-crs}
\end{equation}
Given $(\mathbf u,\mathbf w)\in\cR$, compute 
$(\mathsf A,\mathsf B)\leftarrow\Enc(\mathbf u,\mathbf w)$, express the two
polynomials in the fixed bases as $\mathbf c_A\in \F^{d_U}$ and
$\mathbf c_B\in \F^{d_V}$, and compute two proof halves
\begin{equation}
  A=\mathsf A(\mathbf x)=\mathbf c_A\boldsymbol{\sigma}_1^T,
  \qquad
  B=\mathsf B(\mathbf x)=\mathbf c_B\boldsymbol{\sigma}_2^T.
  \label{eq:generic-fsp-proof}
\end{equation}
The verifier performs any explicit domain check required for
$\mathbf u\in\mathcal X$ and accepts if
\begin{equation}
  AB=\boldsymbol{\sigma}_1M_{\mathbf u}\boldsymbol{\sigma}_2^T.
  \label{eq:generic-fsp-verifier}
\end{equation}
By definition of $M_{\mathbf u}$, the right-hand side equals
$H_{\mathbf u}(\mathbf x)$, so formal factorization gives perfect
completeness.

A split-affine strategy is described by one coefficient row for each CRS
half.  Since $1\in\cU\cap\cV$, these rows define polynomials
$\mathsf A\in\cU$ and $\mathsf B\in\cV$, including arbitrary affine
constants.  If
\[
  \mathsf A(\mathbf X)\mathsf B(\mathbf X)-H_{\mathbf u}(\mathbf X)
\]
is the zero polynomial, $\Dec$ returns a witness.  Otherwise it is a nonzero
polynomial of degree at most $d_{\FSP}$ and vanishes at the sampled point with
probability at most $d_{\FSP}/q$ by Schwartz--Zippel.

Finally, a fixed cross-bilinear zero test on the CRS has formal polynomial
\[
  P_T(\mathbf X)=\mathbf U(\mathbf X)T\mathbf V(\mathbf X)^T
\]
for a setup-independent matrix $T\in \F^{d_U\times d_V}$.  It is either
formally zero, in which case the disclosure event is impossible, or it
vanishes accidentally with probability at most $d_{\FSP}/q$.  The event in
Definition~\ref{def:nilp-security} requires the polynomial to vanish at the
first sampled point, so its probability is at most $d_{\FSP}/q$.
\end{proof}

\section{A factorization span program for signed linear systems}
\label{sec:concrete-fsp}

We now construct an FSP realizing the relation from
Section~\ref{sec:boolean-system}.  The formal
variables are $X,Y,Z$, and their independently sampled setup evaluations will
be denoted by $x,y,z$.

\subsection{Public extendability}

The target-span argument needs a field-valued extension for every public
input, including false Boolean instances.

\begin{definition}[Public extendability]
\label{def:public-extendability}
A homogeneous signed linear system with public--witness partition $(\ell,n)$
is \emph{publicly extendable} if for every $\mathbf u\in\F^\ell$ there is
$\mathbf w^\star\in\F^n$ such that
\[
  \overline{\mathbf a}(\mathbf u,\mathbf w^\star)D=0.
\]
No Booleanity is required of $\mathbf w^\star$.
\end{definition}



Write the rows of $D$ according to the constant, public, and witness
coordinates as
\[
  D=
  \begin{pmatrix}
    \mathbf d_0\\ D_{\mathrm{pub}}\\ D_{\mathrm{wit}}
  \end{pmatrix},
  \qquad
  D_{\mathrm{pub}}=
  \begin{pmatrix}\mathbf d_1\\ \vdots\\ \mathbf d_\ell\end{pmatrix}.
\]

\begin{lemma}[Linear public extension]
\label{lem:linear-public-extension}
If $D$ is publicly extendable, then Gaussian elimination can precompute
$\mathbf b_0,\ldots,\mathbf b_\ell\in\F^n$ satisfying
\[
  \mathbf b_0D_{\mathrm{wit}}=-\mathbf d_0,
  \qquad
  \mathbf b_iD_{\mathrm{wit}}=-\mathbf d_i
  \quad(1\le i\le\ell).
\]
Consequently,
\begin{equation}
  \mathbf w^\star(\mathbf u)
  =\mathbf b_0+\sum_{i=1}^{\ell}u_i\mathbf b_i
  \label{eq:linear-public-extension}
\end{equation}
is an efficiently computable input-linear extension satisfying
$\overline{\mathbf a}(\mathbf u,\mathbf w^\star(\mathbf u))D=0$ for every
$\mathbf u\in\F^\ell$.
\end{lemma}

\begin{proof}
Public extendability at $\mathbf u=\mathbf 0$ gives a vector $\mathbf b_0$
with $\mathbf b_0D_{\mathrm{wit}}=-\mathbf d_0$.  For each standard basis
vector $\mathbf e_i\in\F^\ell$, let $\mathbf c_i$ be an extension.  Then
\[
  \mathbf c_iD_{\mathrm{wit}}=-\mathbf d_0-\mathbf d_i,
\]
so $\mathbf b_i=\mathbf c_i-\mathbf b_0$ has the required property.  These
vectors can be found by Gaussian elimination on the fixed matrix
$D_{\mathrm{wit}}$.  Substitution now proves the final identity.
\end{proof}

\begin{lemma}[The circuit system is publicly extendable]
\label{lem:field-extension}
The matrix $D$ constructed during the Boolean circuit reduction in Section~\ref{sec:boolean-system} is publicly extendable.
\end{lemma}

\begin{proof}
Fix the public coordinates and assign arbitrary field values to the remaining
ordinary wires.  The column belonging to gate $g$ has coefficient $-1$ in its
dedicated residual row $W+g$, and no other gate column uses that row.  The
residual coordinates can therefore be solved independently, one gate equation
at a time.  This is the precise reason a field-valued witness completion always
exists.
\end{proof}

The target-span proof below is most convenient for $n\ge2$.  If necessary,
append unused witness signs whose rows in $D$ are zero; this changes the witness
representation but not the language.

\subsection{Tagged factors for witness signs}

Choose distinct tag points
\[
  \omega_1,\ldots,\omega_N\in\F.
\]
The first $\ell$ tags are reserved for public-input interpolation.  For the
witness tags define
\begin{equation}
  t(Z)=\prod_{j=\ell+1}^{N}(Z-\omega_j),
  \qquad
  t_j(Z)=\frac{t(Z)}{Z-\omega_j}
  \quad(\ell<j\le N).
  \label{eq:fsp-witness-tags}
\end{equation}
Then
\begin{equation}
  t_j(\omega_k)=
  \begin{cases}
    0,&j\ne k,\\
    t'(\omega_j)\neq 0,&j=k,
  \end{cases}
  \qquad \ell<j,k\le N,
  \label{eq:fsp-tag-evaluation}
\end{equation}
where the nonvanishing follows from the distinctness of the tag points.  Also,
$t(\omega_i)\ne0$ for $1\le i\le\ell$.

For a witness assignment
$\mathbf w=(a_{\ell+1},\ldots,a_N)\in \{-1,+1\}^n$, set
\begin{equation}
  F_{\mathbf w}(X,Z)
  =\prod_{j=\ell+1}^{N}\bigl((Z-\omega_j)X+a_j\bigr),
  \qquad
  F_{-\mathbf w}(X,Z)
  =\prod_{j=\ell+1}^{N}\bigl((Z-\omega_j)X-a_j\bigr).
  \label{eq:fsp-selected-factors}
\end{equation}
Their product is independent of the witness:
\begin{equation}
  H(X,Z)
  =\prod_{j=\ell+1}^{N}
    \bigl((Z-\omega_j)X+1\bigr)\bigl((Z-\omega_j)X-1\bigr)
  =\prod_{j=\ell+1}^{N}\bigl((Z-\omega_j)^2X^2-1\bigr),
  \label{eq:fsp-fixed-H}
\end{equation}
so
\begin{equation}
  F_{\mathbf w}(X,Z)F_{-\mathbf w}(X,Z)=H(X,Z).
  \label{eq:fsp-complement-product}
\end{equation}

Writing
\[
  F_{\mathbf w}(X,Z)=\sum_{k=0}^{n}F_{\mathbf w,k}(Z)X^k,
\]
the two leading coefficients are
\begin{equation}
  F_{\mathbf w,n}(Z)=t(Z),
  \qquad
  F_{\mathbf w,n-1}(Z)
  =\sum_{j=\ell+1}^{N}a_jt_j(Z).
  \label{eq:fsp-top-coefficients}
\end{equation}
Hence
\begin{equation}
  \frac{F_{\mathbf w,n-1}(\omega_j)}{t'(\omega_j)}=a_j
  \qquad(\ell<j\le N).
  \label{eq:fsp-recover-witness}
\end{equation}

\begin{lemma}[Elementary factor selection]
\label{lem:factor-selection}
Suppose $F,B\in\F[X,Z]$ satisfy
\[
  F(X,Z)B(X,Z)=H(X,Z),
  \qquad
  \deg_XF=\deg_XB=n.
\]
Then $F$ is a nonzero field scalar times a product of $n$ distinct factors
from
\[
  \bigl\{(Z-\omega_j)X+b:\ell<j\le N,\ b\in\{-1,+1\}\bigr\}.
\]
If its leading coefficient is $v_0t(Z)$, then exactly one factor with each
witness tag occurs and there is a unique
$\mathbf w\in\{-1,+1\}^n$ such that
\[
  F(X,Z)=v_0F_{\mathbf w}(X,Z).
\]
\end{lemma}

\begin{proof}
Temporarily regard $Z$ as symbolic and work with univariate polynomials in
$X$ whose coefficients are rational functions of $Z$.  The $2n$ roots of
$H$ are
\[
  \frac{+1}{Z-\omega_j},
  \qquad
  \frac{-1}{Z-\omega_j},
  \qquad \ell<j\le N,
\]
and are pairwise distinct.  Here $\operatorname{char}(\F)>3$ in
particular implies $+1\ne-1$; together with the distinct tag points this rules
out collisions among the displayed rational functions.  A degree-$n$ factor
therefore selects $n$ of the corresponding linear factors, up to a nonzero
rational scalar $v(Z)$.
The constant coefficient in $X$ of every selected product is $\pm1$.
Since both $F$ and the complementary factor $B=H/F$ have polynomial
coefficients, both $v(Z)$ and $v(Z)^{-1}$ are polynomials; hence $v(Z)$ is a
nonzero field constant.

If $c_j\in\{0,1,2\}$ is the number of selected factors carrying tag $j$, the
leading coefficient of the selected product is
$\prod_j(Z-\omega_j)^{c_j}$.  Equality with $v_0t(Z)$ says that the zero at
each $Z=\omega_j$ has multiplicity one.  Thus $c_j=1$ for every witness tag,
and the linear factors of $F$ determine a unique sign vector $\mathbf w$.
\end{proof}

\subsection{Public-input interpolation and targets}

Assume $\ell\ge1$; a relation with no variable public input can be given one
dummy public coordinate whose domain is the singleton $\{+1\}$.  For
$1\le i\le\ell$, let
$I_i(Z)\in\F^{\le\ell-1}[Z]$ be the Lagrange polynomial on the public tags:
\begin{equation}
  I_i(Z)=
  \prod_{\substack{1\le j\le\ell\\j\ne i}}
  \frac{Z-\omega_j}{\omega_i-\omega_j},
  \qquad
  I_i(\omega_j)=\delta_{ij}.
  \label{eq:fsp-public-lagrange}
\end{equation}
For $\mathbf u\in\F^\ell$, define
\begin{equation}
  J_{\mathbf u}(Z)=\sum_{i=1}^{\ell}u_iI_i(Z).
  \label{eq:fsp-public-interpolant}
\end{equation}
The honest left and right factors are
\begin{equation}
  \mathsf A_{\mathbf u,\mathbf w}(X,Y,Z)
  =\bigl(Y-J_{\mathbf u}(Z)\bigr)F_{\mathbf w}(X,Z),
  \qquad
  \mathsf B_{\mathbf w}(X,Z)=F_{-\mathbf w}(X,Z).
  \label{eq:fsp-honest-factors}
\end{equation}
The target family is
\begin{equation}
  H_{\mathbf u}(X,Y,Z)
  =\bigl(Y-J_{\mathbf u}(Z)\bigr)H(X,Z)
  =YH(X,Z)-\sum_{i=1}^{\ell}u_iI_i(Z)H(X,Z).
  \label{eq:fsp-concrete-target}
\end{equation}
Thus the public input appears linearly, as required by Definition~\ref{def:fsp}.

\subsection{The two fixed spaces}

Define the right triangular space
\begin{equation}
  \cV
  =\Span_{\F}\{X^jZ^k:0\le j\le n,\ 0\le k\le j\}.
  \label{eq:fsp-right-space}
\end{equation}
Equivalently, $\mathsf B=\sum_{j=0}^nB_j(Z)X^j$ belongs to $\cV$ exactly when
$B_j\in\F^{\le j}[Z]$.  Its dimension is
\begin{equation}
  d_V=\dim\cV=\frac{(n+1)(n+2)}2.
  \label{eq:fsp-right-dimension}
\end{equation}

The left space $\cU$ will be a proper subspace of the triangular ambient space
\[
  \Span_{\F}\!\left(
    \{YX^jZ^k:0\le j\le n,\ 0\le k\le j\}
    \cup
    \{X^jZ^k:0\le j\le n,\ 0\le k\le j+\ell-1\}
  \right).
\]
Write a generic left polynomial as
\begin{equation}
  \mathsf A(X,Y,Z)=YF(X,Z)+C(X,Z),
  \label{eq:fsp-generic-left}
\end{equation}
where
\[
  F(X,Z)=\sum_{j=0}^{n}F_j(Z)X^j,
  \qquad
  C(X,Z)=\sum_{j=0}^{n}C_j(Z)X^j.
\]
When $F_n(Z)=v_0t(Z)$, define the encoded scaled assignment
\begin{equation}
  \asg(\mathsf A)
  =\left(
      v_0,
      \left(-\frac{C_n(\omega_i)}{t(\omega_i)}\right)_{i=1}^{\ell},
      \left(\frac{F_{n-1}(\omega_j)}{t'(\omega_j)}\right)_{j=\ell+1}^{N}
    \right)
  \in\F^{N+1}.
  \label{eq:fsp-assignment-map}
\end{equation}
All denominators are fixed nonzero field elements, so this is linear in the
coefficient vector once the leading condition is imposed.

Define $\cU$ to consist of all polynomials $\mathsf A=YF+C$ satisfying
\begin{align}
  F_j&\in\F^{\le j}[Z] &&(0\le j\le n),
  \label{eq:fsp-U-F-degree}\\
  C_j&\in\F^{\le j+\ell-1}[Z] &&(0\le j\le n),
  \label{eq:fsp-U-C-degree}\\
  F_n(Z)&=v_0t(Z) &&\text{for some }v_0\in\F,
  \label{eq:fsp-U-leading}\\
  \asg(\mathsf A)D&=\mathbf 0.
  \label{eq:fsp-U-circuit}
\end{align}
Every condition is homogeneous and linear in the coefficients.  The ambient
left space before the fixed constraints has dimension
\begin{equation}
  \sum_{j=0}^{n}(j+1)+\sum_{j=0}^{n}(j+\ell)
  =(n+1)(N+1).
  \label{eq:fsp-left-ambient-dimension}
\end{equation}
The leading condition has codimension $n$.  After imposing it, the
assignment map is surjective: $v_0$ is chosen through $F_n$, the degree bound
on $C_n$ permits arbitrary values at the $\ell$ public tag points, and the
degree bound on $F_{n-1}$ permits arbitrary values at the $n$ witness tag
points.  These choices use independent coefficient blocks.  The kernel
condition therefore has codimension $\operatorname{rank}(D)$, and
\begin{equation}
  d_U=\dim\cU
  =(n+1)(N+1)-n-\operatorname{rank}(D).
  \label{eq:fsp-left-dimension}
\end{equation}
For the circuit matrix constructed earlier, $\operatorname{rank}(D)=G$.

\begin{remark}[Constants and affine proof maps]
The constant polynomial $1$ lies in both spaces.  For $\cU$, take $F=0$ and
$C=1$, so $v_0=0$ and
$\asg(1)=(0,\ldots,0)\in\F^{N+1}$.  Bases may therefore begin
with $1$, and affine proof maps need no separate CRS coordinates or
multiplicative tags.
\end{remark}

\begin{lemma}[Honest membership]
\label{lem:fsp-honest-membership}
If $(\mathbf u,\mathbf w)\in\cR$, then
\[
  \mathsf A_{\mathbf u,\mathbf w}\in\cU,
  \qquad
  \mathsf B_{\mathbf w}\in\cV.
\]
\end{lemma}

\begin{proof}
The triangular degree bounds follow from multiplying the witness linear
factors.  Equations~\eqref{eq:fsp-top-coefficients} and
\eqref{eq:fsp-public-interpolant} give
\[
  \asg(\mathsf A_{\mathbf u,\mathbf w})
  =\overline{\mathbf a}(\mathbf u,\mathbf w),
\]
so the circuit constraint follows from the relation.  The complementary
factor satisfies the same triangular degree bounds and belongs to $\cV$.
\end{proof}

\subsection{Exact factorization realizes the signed relation}

\begin{theorem}[Realization theorem]
\label{thm:fsp-realization-concrete}
The factor map in \eqref{eq:fsp-honest-factors} and the coefficient decoder
below make the concrete FSP realize $\cR$.
\end{theorem}

\begin{proof}
The encoding direction follows from Lemma~\ref{lem:fsp-honest-membership} and
\eqref{eq:fsp-complement-product}.  Conversely, suppose
\[
  \mathsf A(X,Y,Z)\mathsf B(X,Z)=H_{\mathbf u}(X,Y,Z)
  =\bigl(Y-J_{\mathbf u}(Z)\bigr)H(X,Z)
\]
with $\mathsf A=YF+C\in\cU$ and $\mathsf B\in\cV$.
Comparing coefficients of $Y$ gives
\begin{equation}
  F(X,Z)\mathsf B(X,Z)=H(X,Z).
  \label{eq:fsp-realization-Y}
\end{equation}
The constant coefficient gives
\[
  C(X,Z)\mathsf B(X,Z)
  =-J_{\mathbf u}(Z)H(X,Z)
  =-J_{\mathbf u}(Z)F(X,Z)\mathsf B(X,Z).
\]
Since $\mathsf B\ne0$,
\begin{equation}
  C(X,Z)=-J_{\mathbf u}(Z)F(X,Z).
  \label{eq:fsp-realization-C}
\end{equation}
Both factors in \eqref{eq:fsp-realization-Y} have $X$-degree at most $n$,
whereas $H$ has degree $2n$, so both have degree exactly $n$.  The leading
condition in $\cU$ has $F_n=v_0t$ with $v_0\ne0$.  Lemma~\ref{lem:factor-selection}
therefore gives a unique $\mathbf w\in\{-1,+1\}^n$ with
$F=v_0F_{\mathbf w}$.  Equation~\eqref{eq:fsp-realization-C} then yields
\[
  \asg(\mathsf A)
  =v_0\overline{\mathbf a}(\mathbf u,\mathbf w).
\]
The fixed kernel constraint in $\cU$ implies
$v_0\overline{\mathbf a}(\mathbf u,\mathbf w)D=0$, and $v_0\ne0$ gives
$(\mathbf u,\mathbf w)\in\cR$.

The decoder reads $v_0$ from $F_n=v_0t$ and outputs
\begin{equation}
  a_j=\frac{F_{n-1}(\omega_j)}{v_0t'(\omega_j)}
  \qquad(\ell<j\le N),
  \label{eq:fsp-decoder}
\end{equation}
followed by explicit sign and matrix checks.
\end{proof}

\subsection{Every public target lies in the product span}

\begin{lemma}[Triangular residual support]
\label{lem:triangular-residual}
Let $E(Z)\in\F^{\le n-1}[Z]$ and set
\[
  F_E(X,Z)=t(Z)X^n+E(Z)X^{n-1},
  \qquad
  B_E(X,Z)=t(Z)X^n-E(Z)X^{n-1}.
\]
Every monomial $X^dZ^e$ in
\[
  K_E(X,Z)=H(X,Z)-F_E(X,Z)B_E(X,Z)
\]
satisfies
\[
  0\le d\le2n-2,
  \qquad
  0\le e\le d.
\]
\end{lemma}

\begin{proof}
In $H(X,Z)=\prod_j((Z-\omega_j)^2X^2-1)$, every occurrence of
$X^2$ contributes a polynomial of $Z$-degree at most two.  Hence every
monomial $X^dZ^e$ of $H$ satisfies $e\le d$.  Moreover,
\[
  F_EB_E=t(Z)^2X^{2n}-E(Z)^2X^{2n-2}.
\]
The leading term cancels the $X^{2n}$ term of $H$, and both polynomials have no
$X^{2n-1}$ term.  Finally, $\deg E^2\le2n-2$, so the additional
$X^{2n-2}$ term also satisfies the claimed triangular bound.
\end{proof}

\begin{theorem}[Target-span theorem]
\label{thm:fsp-target-span}
Assume $n\ge2$.  For every $\mathbf u\in\F^\ell$,
\[
  H_{\mathbf u}(X,Y,Z)\in\Span(\cU\cV).
\]
Moreover, for fixed bases $\mathbf U,\mathbf V$ there are efficiently
computable fixed matrices $M_0,M_1,\ldots,M_\ell$ satisfying
\begin{align}
  \mathbf U M_0\mathbf V^T&=YH,
  \label{eq:fsp-M0}\\
  \mathbf U M_i\mathbf V^T&=-I_iH
  \qquad(1\le i\le\ell),
  \label{eq:fsp-Mi}
\end{align}
and hence
\begin{equation}
  M_{\mathbf u}=M_0+\sum_{i=1}^{\ell}u_iM_i
  \label{eq:fsp-Mu}
\end{equation}
represents $H_{\mathbf u}$.
\end{theorem}

\begin{proof}
Fix $\mathbf u\in\F^\ell$ and use
Lemma~\ref{lem:linear-public-extension} to compute
$\mathbf w^\star=\mathbf w^\star(\mathbf u)\in\F^n$ with
$\overline{\mathbf a}(\mathbf u,\mathbf w^\star)D=0$.  Define
\[
  E^\star(Z)=\sum_{j=\ell+1}^{N}a_j^\star t_j(Z),
\]
\[
  F^\star(X,Z)=t(Z)X^n+E^\star(Z)X^{n-1},
  \qquad
  B^\star(X,Z)=t(Z)X^n-E^\star(Z)X^{n-1},
\]
and
\[
  \mathsf A^\star(X,Y,Z)
  =\bigl(Y-J_{\mathbf u}(Z)\bigr)F^\star(X,Z).
\]
The assignment map gives
$\asg(\mathsf A^\star)
 =\overline{\mathbf a}(\mathbf u,\mathbf w^\star)$, so
$\mathsf A^\star\in\cU$; also $B^\star\in\cV$.

By Lemma~\ref{lem:triangular-residual}, the polynomial
\[
  K^\star(X,Z)=H(X,Z)-F^\star(X,Z)B^\star(X,Z)
\]
has only monomials $X^dZ^e$ with $0\le d\le2n-2$ and $0\le e\le d$.
Furthermore,
\[
  \Delta_{\mathbf u}
  :=H_{\mathbf u}-\mathsf A^\star B^\star
  =\bigl(Y-J_{\mathbf u}(Z)\bigr)K^\star(X,Z).
\]
Consequently, a $Y$-monomial in $\Delta_{\mathbf u}$ has the form
$YX^dZ^e$ with $0\le d\le2n-2$ and $0\le e\le d$.
Set $i=\max(0,d-n)$ and $j=d-i$.  Then $0\le i\le n-2$ and
$0\le j\le n$.  Recall that the additional restrictions defining $\cU$
involve only $F_n$, $C_n$, and $F_{n-1}$, hence only the $X^n$ and
  $X^{n-1}$ coefficient blocks.  Since $i\le n-2$, every triangular monomial
  of left $X$-degree $i$ used below belongs to $\cU$.  Set
  $e_L=\max\{0,e-j\}$ and $e_R=e-e_L$.  Then $e_L\le i$ and $e_R\le j$.
\[
  YX^dZ^e=(YX^iZ^{e_L})(X^jZ^{e_R}),
\]
where the left factor is an unrestricted lower $F_i$ monomial in $\cU$ and
the right factor lies in $\cV$.

A $Y$-independent monomial has the form $X^dZ^e$ with the same range of $d$
and $0\le e\le d+\ell-1$.  Set $e_L=\max\{0,e-j\}$ and $e_R=e-e_L$.  Then
$e_L\le i+\ell-1$ and $e_R\le j$.
\[
  X^dZ^e=(X^iZ^{e_L})(X^jZ^{e_R}),
\]
where the left factor is an unrestricted lower $C_i$ monomial in $\cU$.
Thus every monomial of $\Delta_{\mathbf u}$ belongs to
$\Span(\cU\cV)$, proving the first claim.

Setting $\mathbf u=\mathbf 0$ gives $H_{\mathbf 0}=H_0=YH$, and hence a
product-span representation of $YH$.  For $1\le i\le\ell$, let
$\mathbf e_i\in\F^\ell$
denote the $i$th standard basis vector.  By the definition of the target
family,
\[
  H_{\mathbf e_i}=YH-I_iH,
\]
so subtracting a representation of $H_{\mathbf e_i}$ from the representation
of $YH$ gives a representation of $I_iH$.  Negating this representation and
writing all representations in the fixed bases $\mathbf U$ and $\mathbf V$
gives the claimed matrices $M_0,M_1,\ldots,M_\ell$.
\end{proof}

\begin{theorem}[Concrete FSP]
\label{thm:concrete-fsp}
The spaces $\cU,\cV$, the target components
\[
  H_0(X,Y,Z)=YH(X,Z),
  \qquad
  H_i(X,Y,Z)=-I_i(Z)H(X,Z),
\]
and the matrices $M_0,M_1,\ldots,M_\ell$ from
Theorem~\ref{thm:fsp-target-span} form an input-linear FSP realizing
$\cR$.  Its dimensions satisfy
\[
  d_U=(n+1)(N+1)-n-\operatorname{rank}(D),
  \qquad
  d_V=\frac{(n+1)(n+2)}2,
\]
and its degree parameter is bounded by
\begin{equation}
  d_{\FSP}
  \le 4n+\max\{1,\ell-1\}
  \le4N+1.
  \label{eq:concrete-fsp-degree}
\end{equation}
\end{theorem}

\begin{proof}
Realization is Theorem~\ref{thm:fsp-realization-concrete}, and target
representability is Theorem~\ref{thm:fsp-target-span}.  The dimensions were computed above.  The triangular bounds give
\[
  \deg_{\mathrm{tot}}(YF)\le2n+1,
  \qquad
  \deg_{\mathrm{tot}}C\le2n+\ell-1,
  \qquad
  \deg_{\mathrm{tot}}\mathsf B\le2n.
\]
Hence
\[
  \deg_{\mathrm{tot}}(\mathsf A\mathsf B)
  \le\max\{4n+1,4n+\ell-1\}
  =4n+\max\{1,\ell-1\}.
\]
The same bound covers the target components $YH$ and $I_iH$.
\end{proof}

\section{The split NILP and pairing argument}
\label{sec:nilp}

The field-level proof system is now the compiler of
Theorem~\ref{thm:fsp-compiler} applied to Theorem~\ref{thm:concrete-fsp}.
We spell it out to make the split structure and the short verifier explicit.

\subsection{Field-level NILP algorithms}

Choose bases
\[
  \mathbf U(X,Y,Z)=(U_1,\ldots,U_{d_U}),
  \qquad
  \mathbf V(X,Z)=(V_1,\ldots,V_{d_V})
\]
with $U_1=V_1=1$.  Instantiating the syntax above, $\Setup(\cR)$ samples
independently
\[
  x,y,z\leftarrow\F
\]
and computes
\begin{equation}
  \boldsymbol{\sigma}_1=\mathbf U(x,y,z),
  \qquad
  \boldsymbol{\sigma}_2=\mathbf V(x,z),
  \qquad
  \td=(x,y,z).
  \label{eq:concrete-nilp-crs}
\end{equation}

On input $(\mathbf u,\mathbf w)\in\cR$, the algorithm $\ProofMatrix$ samples
$r\leftarrow\F^\times$ and forms
\begin{equation}
  \widetilde{\mathsf A}=r\mathsf A_{\mathbf u,\mathbf w},
  \qquad
  \widetilde{\mathsf B}=r^{-1}\mathsf B_{\mathbf w}.
  \label{eq:concrete-poly-randomization}
\end{equation}
Let $\mathbf c_A\in\F^{d_U}$ and $\mathbf c_B\in\F^{d_V}$ be their unique
coefficient rows in the fixed bases, so
\[
  \widetilde{\mathsf A}=\mathbf c_A\mathbf U^T,
  \qquad
  \widetilde{\mathsf B}=\mathbf c_B\mathbf V^T.
\]
It returns the two proof matrices
\[
  P_1=\mathbf c_A\in\F^{1\times d_U},
  \qquad
  P_2=\mathbf c_B\in\F^{1\times d_V}.
\]
Thus $\ProofMatrix$ depends on $(\mathbf u,\mathbf w,r)$ but not on the
sampled evaluation point $(x,y,z)$ or the CRS values.  The prover algorithm
\[
  \Prove(\cR,\boldsymbol{\sigma}_1,\boldsymbol{\sigma}_2,
         \mathbf u,\mathbf w)
\]
runs $\ProofMatrix(\cR,\mathbf u,\mathbf w)$, applies the resulting matrices
to the corresponding CRS halves, and outputs
\begin{equation}
  \ProofObj=(A,B),
  \qquad
  A=\boldsymbol{\sigma}_1P_1^T
   =\widetilde{\mathsf A}(x,y,z),
  \qquad
  B=\boldsymbol{\sigma}_2P_2^T
   =\widetilde{\mathsf B}(x,z).
  \label{eq:concrete-field-proof}
\end{equation}
For an instance $\mathbf u$, the algorithm $\Test$ returns the single matrix
\[
  T_{\mathbf u}
  =\begin{pmatrix}
     M_{\mathbf u} & 0\\
     0 & -1
  \end{pmatrix},
  \qquad
  M_{\mathbf u}=M_0+\sum_{i=1}^{\ell}u_iM_i,
\]
where the zero blocks have dimensions $d_U\times1$ and $1\times d_V$.
On input $(\boldsymbol{\sigma}_1,\boldsymbol{\sigma}_2,\mathbf u,
\ProofObj)$, the verifier algorithm $\Vfy$ rejects unless every
$u_i\in\{-1,+1\}$ and otherwise accepts precisely when
\begin{equation}
  (\boldsymbol{\sigma}_1,A)T_{\mathbf u}
  (\boldsymbol{\sigma}_2,B)^T=0,
  \qquad\text{equivalently}\qquad
  AB=\boldsymbol{\sigma}_1M_{\mathbf u}\boldsymbol{\sigma}_2^T.
  \label{eq:concrete-field-verifier}
\end{equation}
Each proof matrix has one row, so the proof has split dimensions $(1,1)$; the
verifier evaluates one cross-bilinear equation.

\begin{corollary}[Field-level security]
\label{cor:field-security}
The NILP has perfect completeness, split-affine statistical knowledge
error at most $d_{\FSP}/q$, and disclosure disagreement at most
$d_{\FSP}/q$.
\end{corollary}

\begin{proof}
Apply Theorem~\ref{thm:fsp-compiler}.  The decoder is given explicitly in
\eqref{eq:fsp-decoder}.
\end{proof}

\subsection{Compilation to Type-III pairings}

We use Groth's generic bilinear-group compiler~\cite[Section~2.5]{Groth16} to derive a preprocessing zk-SNARK from the $(1,1)$-split NILP.

Let
\[
  e:\Gone\times\Gtwo\longrightarrow\GT
\]
be a Type-III bilinear map whose three groups have prime order $q$, and take
$\F=\F_q$.  We use additive notation for all three groups.  Fix generators
$G_1\in\Gone$ and $G_2\in\Gtwo$, and let $G_T=e(G_1,G_2)\in\GT$.  Write
$[a]_1=aG_1$, $[a]_2=aG_2$, and $[a]_T=aG_T$, so
$e([a]_1,[b]_2)=[ab]_T$.

Encode the two CRS halves as
\begin{equation}
  \boldsymbol{\sigma}_1=[\mathbf U(x,y,z)]_1,
  \qquad
  \boldsymbol{\sigma}_2=[\mathbf V(x,z)]_2.
  \label{eq:pairing-prover-crs}
\end{equation}
The proof becomes 
\begin{equation}
  \ProofObj=([A]_1,[B]_2)\in\Gone\times\Gtwo,
  \label{eq:pairing-proof}
\end{equation}
which can be computed linearly from the CRS halves using the proof matrices as in (\ref{eq:concrete-field-proof}).

A full verifier can evaluate the right-hand side of
\eqref{eq:concrete-field-verifier} by a pairing-product computation.
To get a fast online verifier, we can precompute 
\begin{equation}
  \boldsymbol{\sigma}_{\mathrm V}=(K_0,K_1,\ldots,K_\ell),
  \label{eq:short-vk}
\end{equation}
where
\begin{equation}
  K_0=[yH(x,z)]_T,
  \qquad
  K_i=[-I_i(z)H(x,z)]_T.
  \label{eq:short-vk-elements}
\end{equation}
These elements can also be computed directly from $(x,y,z)$ during the setup without materializing the representation matrices $M_0,\ldots,M_\ell$.  By
Theorem~\ref{thm:fsp-target-span}, the same elements are derivable from
the prover CRS using those fixed matrices, so this does not change the security guarantees of the proof system.

Verification therefore becomes
\begin{equation}
  e([A]_1,[B]_2)
  \stackrel{?}{=}
  K_0+\sum_{i=1}^{\ell}u_iK_i.
  \label{eq:short-pairing-verifier}
\end{equation}
The online verifier uses one pairing and $\ell$ target-group
additions or subtractions, independent of $N$.

Groth's disclosure-free split-NILP compiler~\cite[Section~2.5]{Groth16} gives
the compiled argument perfect completeness and
statistical knowledge soundness against adversaries making polynomially many
generic bilinear-group operations.  To interpret the statement
asymptotically, consider a family indexed by a security parameter $\lambda$ in
which $N=N(\lambda)$ is polynomially bounded, $1/q(\lambda)$ is negligible,
and $\log q(\lambda)=\operatorname{poly}(\lambda)$.  Since
$d_{\FSP}\le4N+1$, the field-level knowledge and disclosure errors are then
negligible, as required by the compiler theorem.

\subsection{Rerandomization and statistical zero knowledge}
\label{sec:zk}

The polynomial and proof pairs admit the public rerandomization
\begin{equation}
  (\mathsf A,\mathsf B)\longmapsto(r\mathsf A,r^{-1}\mathsf B),
  \qquad
  (A,B)\longmapsto(rA,r^{-1}B),
  \qquad r\leftarrow\F^\times.
  \label{eq:public-rerandomization}
\end{equation}
The same map is applied by scalar multiplication after group compilation.

\begin{theorem}[Statistical composable zero knowledge]
\label{thm:statistical-zk}
With rerandomization, the NILP satisfies Definition~\ref{def:nilp-security}
with statistical composable zero-knowledge error
\[
  \zeta<\frac{2n}{q}.
\]
Its simulated setup samples the real setup distribution conditioned on
$H(x,z)\ne0$, and proof simulation under the resulting CRS is perfect even
given the simulation trapdoor.
\end{theorem}

\begin{proof}
Let
\begin{equation}
  \rho=\Pr_{x,z\leftarrow\F}[H(x,z)=0].
  \label{eq:zk-bad-event}
\end{equation}
For a fixed witness tag $j$, one of its two factors in $H$ vanishes exactly when
$(z-\omega_j)x\in\{-1,+1\}$.  For each $z\ne\omega_j$ there are exactly two
values of $x$, and for $z=\omega_j$ there are none.  A union bound gives
\begin{equation}
  \rho\le\frac{2n(q-1)}{q^2}<\frac{2n}{q}.
  \label{eq:zk-bad-bound}
\end{equation}

The algorithm $\SimSetup$ rejection-samples $x,y,z\leftarrow\F$ until
$H(x,z)\ne0$, constructs the same CRS as $\Setup$, and retains
$\widehat{\td}=(x,y,z)$.  The statistical distance between the uniform
distribution of $(x,y,z)$ and this conditioned distribution is exactly
$\rho$.  Since the CRS is a deterministic function of these values,
\[
  \Delta(\mathsf{CRS},\widehat{\mathsf{CRS}})
  \le \rho<\frac{2n}{q}.
\]

It remains to prove perfect simulation under the simulated CRS.  Let
 $T=(y-J_{\mathbf u}(z))H(x,z)$.  If $y\ne J_{\mathbf u}(z)$, the simulator samples
$A\leftarrow\F^\times$ uniformly and sets $B=T/A$.  This is exactly the honest
rerandomized proof distribution.  If $y=J_{\mathbf u}(z)$, the honest left proof is always $A=0$ and the right proof is uniform in $\F^\times$; the simulator outputs $(0,B)$
for uniform $B\in\F^\times$.  The distribution is therefore identical for every
valid adaptively chosen instance--witness pair, without the simulator using
the witness.  The equality holds pointwise for every simulated trapdoor, so it
continues to hold when the environment is given that trapdoor.  This proves
perfect proof simulation under $\SimSetup$ and the claimed statistical
composable zero-knowledge error.
\end{proof}

\begin{remark}[Efficiency of simulated setup]
The simulated setup is efficient throughout the parameter range of
Theorem~\ref{thm:main}.  If $q>2n$, rejection sampling takes fewer than
$(1-2n/q)^{-1}\le 2n+1$ trials in expectation.  If $n<q\le2n$, the choice
$x=0$ always satisfies $H(x,z)\ne0$, so a trial succeeds with probability at
least $1/q$ and the expected number of trials is at most $q\le2n$.  This
conditioning affects only the simulator: the real setup remains uniform, so
the knowledge and disclosure bounds proved for it are unchanged.  The pairing
compiler preserves perfect simulation under the simulated CRS by encoding the
two field elements in the corresponding source groups.
\end{remark}

\subsection{Efficiency}

\begin{table}[t]
\centering
\begin{tabularx}{0.94\textwidth}{@{}lX@{}}
\toprule
Object & Size or cost \\
\midrule
Nonconstant signed coordinates & $N=W+G$ \\
Witness coordinates & $n=N-\ell$ \\
Left prover CRS & $\dim\cU=(n+1)(N+1)-n-G$ field or $\Gone$ elements \\
Right prover CRS & $\dim\cV=(n+1)(n+2)/2$ field or $\Gtwo$ elements \\
Proof & two field elements; after compilation, one element in each source group \\
Verifier key & $\ell+1$ target-group elements \\
Online verifier & one pairing and $O(\ell)$ additions/subtractions in $\GT$ \\
Field-level knowledge error & at most $d_{\FSP}/q$, where $d_{\FSP}\le4N+1$ \\
Disclosure disagreement & at most $d_{\FSP}/q$ \\
Composable statistical ZK error & $\zeta<2n/q$ after rerandomization \\
\bottomrule
\end{tabularx}
\caption{Concrete costs for the circuit-specialized construction.  The
displayed formula $d_U=(n+1)(N+1)-n-G$ uses the full-rank circuit matrix,
for which $\operatorname{rank}(D)=G$.}
\label{tab:efficiency}
\end{table}

Table~\ref{tab:efficiency} records the costs for a circuit-specialized setup.  As in QAP-based
preprocessing, the prover CRS depends on the fixed circuit and on which wires
are public, but it can be reused for every instance with that public--witness
partition.  

The table counts encoded CRS elements, not a dense representation of every
fixed linear map in the FSP abstraction.  Since $d_U,d_V=O(N^2)$, one dense
$d_U\times d_V$ matrix $M_i$ may contain $O(N^4)$ field elements, and the
entire family may contain $O(N^5)$ when $\ell=O(N)$.  These matrices witness
product-span membership in the abstract FSP and need not be materialized by
the concrete pairing setup.  After sampling $(x,y,z)$, setup evaluates
$H_0(x,y,z),\ldots,H_\ell(x,y,z)$ directly to obtain the target-group verifier
key in~\eqref{eq:short-vk-elements}.

There remains fixed basis preprocessing for $\cU$.  Its ambient coefficient
space has $O(N^2)$ coordinates and its defining system has $O(N)$ independent
linear constraints.  A dense systematic row reduction therefore gives a
conservative $O(N^4)$-time and $O(N^3)$-working-memory preprocessing bound.
The resulting systematic basis has $O(N^3)$ potentially nonzero field
coefficients, while conversion of a valid polynomial coefficient table to
basis coordinates consists simply of reading its free coordinates.  These
fixed representation costs are separate from the $O(N^2)$ encoded prover CRS
size displayed in Table~\ref{tab:efficiency}.


The right prover CRS has exactly
\[
  d_V=\frac{(n+1)(n+2)}2
\]
source-group elements.  For the circuit matrix, $\operatorname{rank}(D)=G$,
so the left prover CRS has exactly
\[
  d_U=(n+1)(N+1)-n-G
\]
elements.

Each selected or complementary factor polynomial has
$d_V=(n+1)(n+2)/2$ triangular coefficients.  The polynomial
$C=-J_{\mathbf u}F_{\mathbf w}$ has
$\sum_{j=0}^{n}(j+\ell)$ coefficient slots.  After multiplying $k$ tagged
factors, the triangular coefficient table has $\Theta(k^2)$ entries.
Multiplication by one further three-monomial factor therefore costs $O(k^2)$ field
operations, and sequential multiplication costs
$\sum_{k=1}^nO(k^2)=O(n^3)$.  A balanced product tree using standard fast
dense bivariate multiplication reduces this to $\widetilde O(n^2)$ field
operations, quasi-linear in the $\Theta(n^2)$ output table, up to the usual
polylogarithmic factors.

The right space uses its monomial basis.  For the left space, setup uses the
systematic row-echelon basis described above.  A valid polynomial coefficient
table is converted to basis coordinates by reading its free coordinates, not
by solving an instance-dependent linear system.  The compiled prover performs
one multi-scalar multiplication of size $d_U$ in $\Gone$ and one of size
$d_V$ in $\Gtwo$.

Verification checks the public signs, evaluates one pairing, performs
$\ell$ additions or subtractions in $\GT$, and tests one equality.  There
is no online verifier dependence on $N$ when the target group elements
$K_0,K_1,\ldots,K_\ell$ are precomputed.

\section{Related work}

\subsection{Boolean constraints, span programs, and signed linear systems}

Groth, Ostrovsky, and Sahai observed (Section~4.2, Lemma~3) that, for bits
$b_0,b_1,b_2\in\{0,1\}$,
\[
  b_0+b_1+2b_2-2\in\{0,1\}
  \quad\Longleftrightarrow\quad
  b_2=b_0\mathbin{\NAND}b_1,
\]
and used homomorphic commitments to prove that the residual is
Boolean~\cite{GOS06}.
The signed NAND constraint in Section~\ref{sec:boolean-system} is this
observation after the change
of variables $a=2b-1$ and the introduction of an explicit residual sign.
There is also a half-adder interpretation: for a valid NAND gate,
$1-b_2=b_0b_1$ is the carry and
$b_0+b_1+2b_2-2=b_0\oplus b_1$ is the sum.

Danezis, Fournet, Groth, and Kohlweiss~\cite{DFGK14} characterize all
nontrivial fan-in-two Boolean gates through affine maps whose values are
restricted to signed coordinates.  The system $\overline{\mathbf a}D=0$ is a
homogeneous residual-variable form of a similar affine-map characterization.
They reduce the relation to a square span program (SSP), a notion defined in the
same paper, and use pairings to obtain a SNARK.  Our construction instead uses
a factorization span program.

Karchmer and Wigderson introduced span programs as a linear-algebraic model of
Boolean computation~\cite{KW93}.  In a standard span program, the input
assignment determines which labelled vectors are available and acceptance asks
whether a fixed target belongs to their span; the coefficients of the spanning
combination are otherwise arbitrary field elements.  The signed linear system
used here has the complementary form: the matrix $D$ is fixed, while the
assignment is itself the coefficient vector, constrained coordinatewise to
$\{-1,+1\}$, and is required to lie in the left kernel of $D$.

Quadratic span programs (QSPs) and quadratic arithmetic programs (QAPs) provide the main
precedent for relation-specific preprocessing and for compressing many
constraints into one polynomial condition~\cite{GGPR13}.  Square arithmetic
programs (SAPs) exploit the corresponding one-polynomial symmetry in the arithmetic
setting~\cite{GM17}.  Factorization span programs instead distinguish
membership in a product span from the existence of a single exact
factorization.

Lipmaa gives a unified framework for non-universal SNARKs encompassing QAP,
SAP, SSP, and QSP constructions and relates several of their underlying
languages~\cite{Lipmaa22}.  FSP acceptance has a different defining
condition: the public target must have a decomposable preimage, rather than
merely satisfying the quadratic or square-program condition of those
frameworks.

Polynomial products have also appeared before as proof constraints.  Neff's
shuffle argument uses the invariance of a monic polynomial under permutation
of its roots to certify multiset equality~\cite{Neff01}.  Our use of products
is different: the contribution is the complementary tagged-factor FSP and its
split NILP, not the general use of polynomial factorization in zero knowledge.

The conscientious span programs of Gennaro, Gentry, Parno, and Raykova ensure
that every input to a gate checker is actually used, so that wire values can be
extracted consistently~\cite{GGPR13}.  Lipmaa subsequently gave a modular
circuit checker that combines small standard span programs for individual
gates with a high-distance linear error-correcting code enforcing consistency
between wire occurrences~\cite{Lipmaa13}.  In the present formulation, wire
coordinates are shared globally and every gate has a dedicated residual
coordinate.  Gate correctness and wire consistency are thereby expressed by
one sparse homogeneous linear system, without a separate wire checker.

There is also a direct partition interpretation of this system.  If
$\mathbf d_0,\ldots,\mathbf d_N\in\F^G$ are the rows of $D$, then
$\overline{\mathbf a}D=0$ is the signed vector-partition condition
\[
  \sum_{i=0}^{N} a_i\mathbf d_i=0,
  \qquad a_i\in\{-1,+1\}.
\]
Equivalently, the rows assigned sign $+1$ and those assigned sign $-1$ have
equal sums.  Under $x_i=(a_i+1)/2$, this becomes the binary linear-feasibility
equation $\mathbf xD=\frac12\mathbf 1D$, equivalently a multidimensional
subset-sum equation.  Fauzi, Lipmaa, and Zhang constructed a NIZK
argument for the scalar Set-Partition relation
$\sum_i b_i s_i=0$ with $b_i\in\{-1,+1\}$~\cite{FLZ13}.  Our relation is a
sparse vector-valued version with a distinguished constant coordinate and
partially fixed public signs; the proof mechanism is different.

\subsection{From LIPs to split NILPs}

Bitansky, Chiesa, Ishai, Ostrovsky, and Paneth introduced linear interactive
proofs to isolate the information-theoretic linear-prover object behind
several preprocessing SNARKs~\cite{BCI13}.  Groth recast their particular
two-move algebraic input-oblivious LIP as a NILP, placing the verifier's first
message in setup and writing the affine prover strategy as a proof
matrix~\cite{Groth16}.  For that restricted LIP notion, NILP is therefore
largely a change of presentation, but it is not synonymous with LIP in
general.

Dot-product proofs provide another information-theoretic linear proof object:
the verifier makes one fully linear query jointly to the instance and proof.
Their generic-group applications yield very short interactive arguments, but
do not give a publicly verifiable noninteractive Type-III argument with one
element from each source group~\cite{DPP24}.

The split-NILP definition is a substantive refinement for Type-III pairings:
the CRS and proof are partitioned into left and right parts, each proof part
may depend only on its corresponding setup half, and verification is
cross-bilinear.  Two-field-element NILPs already follow from SSPs or
squaring-only circuits.  The open issue is whether the responses can be
separated as one element from each Type-III side.  The contribution is
therefore not the first two-field-element NILP, but an FSP and
resulting disclosure-free NILP with split dimensions exactly $(1,1)$.

The split-NILP perspective has also been used explicitly after Groth.
Groth and Maller formulate pairing-based simulation-extractable arguments in
terms of split NILPs~\cite{GM17}, while Campanelli, Fiore, and Querol restate
the split-NILP compiler and use variants of Groth's NILP to obtain modular
commit-and-prove SNARKs~\cite{LegoSNARK19}.  These works retain at least three
proof elements in the Type-III setting and do not give split dimensions
$(1,1)$.

\subsection{Short pairing-based arguments}

The progression toward short pairing-based arguments includes Groth's initial
constant-size construction, Lipmaa's refinement, the seven-element QSP
argument, the four-element SSP argument, and Groth's three-element
argument~\cite{Groth10,Lipmaa12,GGPR13,DFGK14,Groth16}.  The last remains the
principal Type-III pairing baseline without a random oracle.

Two proof elements were already attainable from squaring-only arithmetizations
using symmetric Type-I pairings~\cite{Groth16}.  This does not yield one
element in each of two distinct source groups and, in standard pairing-friendly
instantiations, the larger symmetric source groups typically make the encoded
proof longer at comparable security.  The question left by Groth is
specifically whether two elements suffice over asymmetric Type-III pairings.

Groth and Maller prove that any pairing-based simulation-extractable SNARK or
signature of knowledge in their model requires at least three group elements
and two verification equations~\cite{GM17}.  Their lower bound concerns the
stronger simulation-extractable notion and therefore does not rule out the
ordinary knowledge-sound $(1,1)$ construction given here.

Among publicly verifiable systems using a random oracle, Polymath uses KZG
commitments for SAPs and, after Fiat--Shamir, has a proof consisting of three
$\Gone$ elements and one field element, with no $\Gtwo$ proof
element (see Table~1 of~\cite{Lipmaa24}); the construction uses KZG
commitments~\cite{KZG10}.  PARI sends two $\Gone$ elements and two field
elements; its noninteractive form also uses the random-oracle
model~\cite{DMS26}.  A later extension adds zero knowledge without changing
this basic proof shape~\cite{zkPARI26}.

Arnon, Dujmovi\'c, and Yogev give a public SNARG with exactly two $\Gone$
elements and no additional bits in the combined generic bilinear-group and
random-oracle model~\cite{ADY26}.  This is the closest result in exact
group-element count.  Because both elements are in $\Gone$, their proof can
also be shorter in encoded bits on common Type-III curves.  Their theorem is
in a different model: the present construction addresses the plain split-NILP
question without a random oracle and has perfect completeness.

\subsection{Short arguments in other models}

Shorter designated verifier arguments exist.  Barta,
Ishai, Ostrovsky, and Wu give a designated-verifier SNARG in the pairing-free
generic-group model with two proof elements and inverse-polynomial soundness
error~\cite{BIOW20}.  Arnon, Dujmovi\'c, and Ishai later obtain negligible
soundness with one group element and additional security-parameter-length
bits, again for a designated verifier~\cite{ADI25}.  These constructions rely
on a secret verification key and do not address the publicly verifiable
Type-III problem.

Sahai and Waters give a qualitatively different route to short arguments from
indistinguishability obfuscation and puncturable pseudorandom
functions~\cite{SW14}.  This is a separate feasibility route based on
incomparable assumptions and proof techniques.


%Finally, the generic-group lineage should credit both Nechaev and Shoup.
%Nechaev proved a generic lower bound for deterministic discrete-logarithm
%algorithms in 1994~\cite{Nec94}; Shoup subsequently gave the standard
%randomized generic-group formulation and matching lower bounds in
%1997~\cite{Sho97}.  Maurer later gave an abstract register-based model of
%generic computation~\cite{Maurer05}, while Boneh, Boyen, and Goh developed the
%polynomial method commonly used for generic bilinear-group analyses~\cite{BBG05}.
%The pairing-level knowledge claim here uses the random-encoding formulation and
%follows Groth's generic bilinear-group treatment; it should not be read as a
%standard-model security claim.


\section{Conclusion}

\begin{theorem}[Field-level main theorem]
\label{thm:main}
Let $\ell\ge1$, $n\ge2$, and $N=n+\ell$.  Let $\F$ be a finite
field of size $q>N$ and characteristic greater than $3$, and let
$D\in\F^{(N+1)\times G}$.  Suppose the signed homogeneous system defined by
$D$ is publicly extendable in the sense of
Definition~\ref{def:public-extendability}.  The construction gives an
input-linear FSP realizing $\cR$ and, by random evaluation, a $(1,1)$-split
NILP with:
\begin{enumerate}[label=(\roman*)]
  \item perfect completeness;
  \item split-affine statistical knowledge error at most
  $d_{\FSP}/q$, where
  $d_{\FSP}\le4n+\max\{1,\ell-1\}\le4N+1$;
  \item disclosure disagreement at most $d_{\FSP}/q$;
  \item statistical composable zero knowledge with error $\zeta<2n/q$ under
  multiplicative rerandomization;
  \item prover CRS size
  \[
    d_U+d_V
    =(n+1)(N+1)-n-\operatorname{rank}(D)
      +\frac{(n+1)(n+2)}2,
  \]
  one field response from each split side, and an input-linear target family
  whose $\ell+1$ component polynomials admit fixed representation matrices.
  The concrete pairing setup need not materialize these matrices.
\end{enumerate}
\end{theorem}

Under the minimal condition $q>N$, the displayed algebraic error bounds may
be vacuous for small fields.  They are intended to become negligible in the
asymptotic instantiation below.

\begin{corollary}[Boolean circuits]
Let $\cC$ be a Boolean circuit with $G$ gates, $\ell\ge1$ public inputs, and
$W$ input and internal wires after eliminating the accepting output wire, and
suppose $n:=G+W-\ell\ge2$.  Over any field satisfying the hypotheses of
Theorem~\ref{thm:main}, the construction gives a disclosure-free $(1,1)$-split
NILP for $\cR_{\cC}$ with the same properties, where $N=G+W$ and
$\operatorname{rank}(D)=G$.
\end{corollary}

\begin{corollary}[Type-III pairing compilation]
\label{cor:pairing-compilation}
Assume in addition that $q$ is prime and that Type-III bilinear groups of
order $q$ are available.  For a family indexed by a security parameter
$\lambda$, suppose also that $N=N(\lambda)$ is polynomially bounded and
$1/q(\lambda)$ is negligible, while
$\log q(\lambda)=\operatorname{poly}(\lambda)$.  Encoding the two field-level
CRS halves and proof responses in the corresponding source groups gives a
pairing argument whose proof consists of one $\Gone$ element and one $\Gtwo$ element.
The derived verifier key contains $\ell+1$ target-group elements, and online
verification uses one pairing and a linear combination in the public input.
The argument has perfect completeness, preserves the field-level statistical
composable zero knowledge with negligible error, and, by Groth's split-NILP
compiler, is statistically knowledge sound against adversaries making
polynomially many generic bilinear-group operations.
\end{corollary}

The main result closes the split-dimension gap at the NILP level.  Through
Groth's compiler, it also answers the corresponding generic Type-III pairing
question.

\section*{AI declaration}
AI was used extensively in this research. Prompting GPT 5.6 Sol Pro (August 2026) with the open question yielded a split NILP for an exact knapsack problem leveraging the core ideas in the FSP. Further prompting yielded a split NILP for Boolean circuit satisfiability, and subsequently reduced the verification cost from $O(N^2)$ to $O(\ell)$. It is still an open question whether it is possible to keep the advantages of the proof system but reduce the CRS size to $O(N)$. Sol Pro found a complicated way to reduce the size to $O(N^2/\log N)$ but did not make further progress.

\begin{thebibliography}{99}

\bibitem[ADI25]{ADI25}
Gal Arnon, Jesko Dujmovi\'c, and Yuval Ishai.
\newblock Designated-Verifier SNARGs with One Group Element.
\newblock In \emph{CRYPTO 2025, Part VII}, LNCS 16006, pages 192--224, 2025.
\newblock \url{https://doi.org/10.1007/978-3-032-01907-3_7}.

\bibitem[ADY26]{ADY26}
Gal Arnon, Jesko Dujmovi\'c, and Eylon Yogev.
\newblock Pairing-Based SNARGs with Two Group Elements.
\newblock In \emph{CRYPTO 2026, Part IX}, LNCS 16808, pages 125--156, 2026.
\newblock Cryptology ePrint Archive, Paper 2025/2160.
\newblock \url{https://eprint.iacr.org/2025/2160}.

\bibitem[BIOW20]{BIOW20}
Ohad Barta, Yuval Ishai, Rafail Ostrovsky, and David J. Wu.
\newblock On Succinct Arguments and Witness Encryption from Groups.
\newblock In \emph{CRYPTO 2020, Part I}, LNCS 12170, pages 776--806, 2020.
\newblock \url{https://eprint.iacr.org/2020/1319}.

\bibitem[BCI+13]{BCI13}
Nir Bitansky, Alessandro Chiesa, Yuval Ishai, Rafail Ostrovsky, and
Omer Paneth.
\newblock Succinct Non-Interactive Arguments via Linear Interactive Proofs.
\newblock In \emph{TCC 2013}, LNCS 7785, pages 315--333, 2013.
\newblock \url{https://eprint.iacr.org/2012/718}.

\bibitem[BHI+24]{DPP24}
Nir Bitansky, Prahladh Harsha, Yuval Ishai, Ron D. Rothblum, and David J. Wu.
\newblock Dot-Product Proofs and Their Applications.
\newblock In \emph{FOCS 2024}, pages 806--825, 2024.
\newblock \url{https://doi.org/10.1109/FOCS61266.2024.00057}.

\bibitem[CFQ19]{LegoSNARK19}
Matteo Campanelli, Dario Fiore, and Ana\"{\i}s Querol.
\newblock LegoSNARK: Modular Design and Composition of Succinct
Zero-Knowledge Proofs.
\newblock In \emph{ACM CCS 2019}, pages 2075--2092, 2019.
\newblock \url{https://doi.org/10.1145/3319535.3339820}.

\bibitem[DFGK14]{DFGK14}
George Danezis, C\'{e}dric Fournet, Jens Groth, and Markulf Kohlweiss.
\newblock Square Span Programs with Applications to Succinct NIZK Arguments.
\newblock In \emph{ASIACRYPT 2014}, LNCS 8873, pages 532--550, 2014.
\newblock \url{https://eprint.iacr.org/2014/718}.

\bibitem[DMS26]{DMS26}
Michel Dellepere, Pratyush Mishra, and Alireza Shirzad.
\newblock GARUDA and PARI: Faster and Smaller SNARKs via Equifficient
Polynomial Commitments.
\newblock In \emph{USENIX Security 2026}, 2026.
\newblock Cryptology ePrint Archive, Paper 2024/1245.
\newblock \url{https://eprint.iacr.org/2024/1245}.

\bibitem[FLZ13]{FLZ13}
Prastudy Fauzi, Helger Lipmaa, and Bingsheng Zhang.
\newblock Efficient Modular NIZK Arguments from Shift and Product.
\newblock In \emph{CANS 2013}, LNCS 8257, pages 92--121, 2013.
\newblock Cryptology ePrint Archive, Paper 2012/548.
\newblock \url{https://eprint.iacr.org/2012/548}.

\bibitem[GGPR13]{GGPR13}
Rosario Gennaro, Craig Gentry, Bryan Parno, and Mariana Raykova.
\newblock Quadratic Span Programs and Succinct NIZKs without PCPs.
\newblock In \emph{EUROCRYPT 2013}, LNCS 7881, pages 626--645, 2013.
\newblock \url{https://eprint.iacr.org/2012/215}.

\bibitem[Gro06]{Groth06CZK}
Jens Groth.
\newblock Simulation-Sound NIZK Proofs for a Practical Language and Constant
Size Group Signatures.
\newblock In \emph{ASIACRYPT 2006}, LNCS 4284, pages 444--459, 2006.
\newblock \url{https://doi.org/10.1007/11935230_29}.

\bibitem[Gro10]{Groth10}
Jens Groth.
\newblock Short Pairing-Based Non-Interactive Zero-Knowledge Arguments.
\newblock In \emph{ASIACRYPT 2010}, LNCS 6477, pages 321--340, 2010.
\newblock \url{https://eprint.iacr.org/2009/390}.

\bibitem[Gro16]{Groth16}
Jens Groth.
\newblock On the Size of Pairing-based Non-interactive Arguments.
\newblock In \emph{EUROCRYPT 2016, Part II}, LNCS 9666, pages 305--326,
2016.
\newblock Cryptology ePrint Archive, Paper 2016/260.
\newblock \url{https://eprint.iacr.org/2016/260}.

\bibitem[GM17]{GM17}
Jens Groth and Mary Maller.
\newblock Snarky Signatures: Minimal Signatures of Knowledge from
Simulation-Extractable SNARKs.
\newblock In \emph{CRYPTO 2017, Part II}, LNCS 10402, pages 581--612, 2017.
\newblock \url{https://eprint.iacr.org/2017/540}.

\bibitem[GOS06]{GOS06}
Jens Groth, Rafail Ostrovsky, and Amit Sahai.
\newblock Perfect Non-interactive Zero Knowledge for NP.
\newblock In \emph{EUROCRYPT 2006}, LNCS 4004, pages 339--358, 2006.
\newblock \url{https://eprint.iacr.org/2005/290}.

\bibitem[KW93]{KW93}
Mauricio Karchmer and Avi Wigderson.
\newblock On Span Programs.
\newblock In \emph{8th Annual Structure in Complexity Theory Conference},
pages 102--111, 1993.
\newblock \url{https://doi.org/10.1109/SCT.1993.336536}.

\bibitem[KZG10]{KZG10}
Aniket Kate, Gregory M. Zaverucha, and Ian Goldberg.
\newblock Constant-Size Commitments to Polynomials and Their Applications.
\newblock In \emph{ASIACRYPT 2010}, LNCS 6477, pages 177--194, 2010.
\newblock \url{https://doi.org/10.1007/978-3-642-17373-8_11}.

\bibitem[KPS26]{zkPARI26}
Shubham Khurana, Sahadeo Padhye, and Rajeev Anand Sahu.
\newblock Zero-Knowledge Extension of PARI.
\newblock \emph{IACR Communications in Cryptology}, 3(2), 2026.
\newblock \url{https://doi.org/10.62056/ahjb0l2hd}.

\bibitem[Lip12]{Lipmaa12}
Helger Lipmaa.
\newblock Progression-Free Sets and Sublinear Pairing-Based Non-Interactive
Zero-Knowledge Arguments.
\newblock In \emph{TCC 2012}, LNCS 7194, pages 169--189, 2012.
\newblock \url{https://doi.org/10.1007/978-3-642-28914-9_10}.

\bibitem[Lip13]{Lipmaa13}
Helger Lipmaa.
\newblock Succinct Non-Interactive Zero Knowledge Arguments from Span
Programs and Linear Error-Correcting Codes.
\newblock In \emph{ASIACRYPT 2013, Part I}, LNCS 8269, pages 41--60, 2013.
\newblock \url{https://doi.org/10.1007/978-3-642-42033-7_3}.

\bibitem[Lip22]{Lipmaa22}
Helger Lipmaa.
\newblock A Unified Framework for Non-universal SNARKs.
\newblock In \emph{PKC 2022, Part I}, LNCS 13177, pages 553--583, 2022.
\newblock \url{https://doi.org/10.1007/978-3-030-97121-2_20}.

\bibitem[Lip24]{Lipmaa24}
Helger Lipmaa.
\newblock Polymath: Groth16 Is Not the Limit.
\newblock In \emph{CRYPTO 2024, Part X}, LNCS 14929, pages 170--206, 2024.
\newblock \url{https://eprint.iacr.org/2024/916}.

\bibitem[Nef01]{Neff01}
C. Andrew Neff.
\newblock A Verifiable Secret Shuffle and Its Application to E-Voting.
\newblock In \emph{ACM CCS 2001}, pages 116--125, 2001.
\newblock \url{https://doi.org/10.1145/501983.502000}.

\bibitem[SW14]{SW14}
Amit Sahai and Brent Waters.
\newblock How to Use Indistinguishability Obfuscation: Deniable Encryption,
and More.
\newblock In \emph{STOC 2014}, pages 475--484, 2014.
\newblock \url{https://eprint.iacr.org/2013/454}.

\bibitem[Val76]{Valiant76}
Leslie G. Valiant.
\newblock Universal Circuits (Preliminary Report).
\newblock In \emph{STOC 1976}, pages 196--203, 1976.
\newblock \url{https://doi.org/10.1145/800113.803649}.

\end{thebibliography}

\end{document}
